% !TEX root = ./index.tex

\maketitle
\newpage

\section*{Answers}

\begin{enumerate}

    \item Let $G$ be a simple connected graph.
        For any two vertices $u, v$ of $G$, define their \emph{distance} as the length of the shortest path connecting them.
        Then, the diameter and radius of the graph are defined as follows:

        \smallskip

        \begin{definition} \label{def:graph-diameter}
            The \emph{diameter} of the graph is
            \[ \mathrm{diam}(G) := \max_{u, v \in V(G)} d(u, v). \]
        \end{definition}

        \begin{definition} \label{def:eccentricity}
            The \emph{eccentricity} of a vertex $v \in V(G)$ is
            \[ e(v) := \max_{u \in V(G)} d(u, v). \]
        \end{definition}

        \begin{definition} \label{def:graph-radius}
            The \emph{radius} of the graph is
            \[ r(G) := \min_{v \in V(G)} e(v). \]
        \end{definition}

        \bigskip

        \begin{enumerate}

            \item There are six non-isomorphic simple connected graphs with four vertices.
                Their diameters and radii are listed in \cref{fig:graphs}.

                \begin{figure}[ht]
                    \centering
                    \vspace{10pt}
                    \includegraphics[width=0.8\textwidth]{figures/graphs-with-four-vertices.tikz}
                    \caption{The diameter $D$ and radius $r$ of each simple connected graph with four vertices.}
                    \label{fig:graphs}
                \end{figure}

                \bigskip

            \item
                \begin{proposition} \label{prop:1-b}
                    $r(G) \le \mathrm{diam}(G) \le 2r(G)$.
                \end{proposition}

                \medskip

                \begin{proof}
                    The left-hand inequality follows from the fact that
                    \[ \mathrm{diam}(G) = \max_{v \in V(G)} \left[\max_{u \in V(G)} d(u, v)\right] = \max_{v \in V(G)} e(v). \]

                    From this, we have
                    \[ r(G) = \min_{v \in V(G)} e(v) \le \max_{v \in V(G)} e(v) = \mathrm{diam}(G). \]

                    \medskip

                    Now to prove the right-hand inequality.

                    For any sequence of vertices $A = (a_0, a_1, \ldots, a_n)$ such that $a_i \in V(G)$, we define $P(A)$ to be the length of the shortest walk connecting each vertex in order.

                    It can be shown that for any $0 \le i \le n$,
                    \begin{equation} \label{eq:walks}
                        P((a_0, \ldots a_i)) + P((a_i, \ldots, a_n)) = P(A).
                    \end{equation}

                    \smallskip

                    Now, let $p, q$ be two vertices in $V(G)$ such that $d(p, q) = \mathrm{diam}(G)$, and let $s$ be some vertex in $V(G)$ such that $e(s) = r(G)$.
                    Note that $s$ may be equal to $p$ or $q$.

                    Define the sequence $X = (p, s, q)$. From \cref{eq:walks}, we have
                    \[ P(p, s) + P(s, q) = P(X). \]
                    The shortest walk between two vertices $u, v$ has length $d(u, v)$, and by definition, $d(u, v) \le e(v)$.
                    Hence,
                    \begin{align*}
                        P(X) &= d(p, s) + d(s, q), \\
                        &\le e(s) + e(s), \\
                        &= 2r(G).
                    \end{align*}

                    Since the shortest path from $p$ to $q$ cannot be longer than the shortest walk from $p$ to $s$ then to $q$, we have $d(p, q) \le P(X)$.

                    Combining these inequalities gives
                    \[ \mathrm{diam}(G) = d(p, q) \le P(X) \le 2r(G), \]
                    which is the desired result.
                \end{proof}

                \bigskip

            \item
                \begin{proposition} \label{prop:1-c}
                    For any choice of $r, D \in \mathbb{N}$ such that $r \le D \le 2r$, there exists a graph $G$ such that $r(G) = r$ and $\mathrm{diam}(G) = D$.
                \end{proposition}

                \medskip

                \begin{proof}
                    Consider the graph $G$ with vertices
                    \[ V = \{v_0, v_1, \ldots, v_{2r}\} \]
                    and edges
                    \[ E = \{\{v_0, v_1\}, \{v_1, v_2\}, \ldots, \{v_{2r - 1}, v_{2r}\}\}. \]
                    This graph resembles a chain of vertices.

                    \begin{figure}[ht]
                        \centering
                        \vspace{10pt}
                        \includegraphics[width=0.6\textwidth]{figures/graph-with-given-radius.tikz}
                        \caption{The graph $G(V, E)$ with $r = 4$.}
                        \label{fig:graph-with-given-radius}
                    \end{figure}

                    The radius of this graph is equal to $e(v_r)$, since any other vertex is further away than $v_r$ from either $v_0$ or $v_{2r}$.
                    Since the distance from $v_r$ to an endpoint of the graph is $r$, the radius is also $r$.

                    \smallskip

                    To force the diameter of the graph to be equal to $D$, we shall add the edge $\{v_k, v_{2r}\}$ to $E$, where $k = 2(D - r)$.

                    \begin{figure}[ht]
                        \centering
                        \vspace{10pt}
                        \includegraphics[width=0.6\textwidth]{figures/graph-with-given-diameter-and-radius.tikz}
                        \caption{The new graph $G(V, E)$ with $r = 4$ and $D = 6$.}
                        \label{fig:graph-with-given-diameter-and-radius}
                    \end{figure}

                    From \cref{fig:graph-with-given-diameter-and-radius}, the diameter of the new graph is equal to the eccentricity of the remaining endpoint $v_0$.
                    We will show that this is equal to $D$.

                    Note that since $D \le 2r$,
                    \begin{align*}
                        k &= 2(D - r), \\
                        &= 2D - 2r, \\
                        &\le 2D - D, \\
                        &= D.
                    \end{align*}

                    In the case where $k = D$, the extra edge does not affect the shortest path from $v_0$ to $v_D$, so $d(v_0, v_D) = D$.

                    Otherwise, the edge $\{v_k, v_{2r}\}$ connects 'across' $v_D$ so that there are two paths between $v_0$ and $v_D$: one with the edge and one without it.

                    From the graph, the length of the path without the edge is simply $D$.

                    The length of the path with the edge is equal to
                    \begin{align*}
                        d(v_0, v_k) + 1 + d(2r, v_D) &= k + 1 + (2r - D) \\
                        &= 2(D - r) + 1 + (2r - D) \\
                        &= D + 1.
                    \end{align*}

                    Thus, the shortest path from $v_0$ to $v_k$ has length $D$.
                    This means that $d(v_0, v_D)$ is always $D$.

                    Changing $v_D$ to any other vertex $v_i$ with $k \le i \le 2r$ must shorten one of the two paths leading to the vertex, hence the value of $d(v_0, v_i)$ can never be greater than $d(v_0, v_D)$.
                    Thus,
                    \[ \mathrm{diam}(G) = e(v_0) = d(v_0, v_D) = D. \]

                    The constructed graph $G$ therefore satisfies the given criteria.
                \end{proof}

        \end{enumerate}

        \bigskip

    \item For any natural number $n$, denote by $T(n)$ the number of non-congruent triangles with sides of integer lengths and perimeter $n$.

        \bigskip

        \begin{enumerate}

            \item Consider the natural numbers $a, b, c$ such that $a \le b \le c$.
                These numbers describe the lengths of a triangle if and only if they obey the triangle inequality $a + b > c$.
                Adding $c$ to both sides gives
                \begin{spreadlines}{\baselineskip}
                    \begin{align*}
                        && a + b + c &> 2c, && \\
                        \implies && \frac{a + b + c}{2} &> c.
                    \end{align*}
                \end{spreadlines}

                The value of $T(n)$ is thus the number of distinct sets $\{a, b, c\}$ of natural numbers such that $a + b + c = n$ and $\max(\{a, b, c\}) < \frac{n}{2}$.
                These sets must be distinct because if they are not, then the triangles they describe are congruent.

                Let $L_n$ denote the set of all possible sets that fulfull this condition for some $n$.
                Calculating $L_n$ for each $1 \le n \le 10$:
                \begin{spreadlines}{0.5\baselineskip}
                    \begin{flalign*}
                        L_1 &= \varnothing, \\
                        L_2 &= \varnothing, \\
                        L_3 &= \{\{1, 1, 1\}\}, \\
                        L_4 &= \{\{1, 1, 2\}\}, \\
                        L_5 &= \{\{1, 2, 2\}\}, \\
                        L_6 &= \{\{2, 2, 2\}\}, \\
                        L_7 &= \{\{1, 3, 3\}, \{2, 3, 3\}\}, \\
                        L_8 &= \{\{2, 3, 3\}\}, \\
                        L_9 &= \{\{1, 4, 4\}, \{2, 3, 4\}, \{3, 3, 3\}\}, \\
                        L_{10} &= \{\{2, 4, 4\}, \{3, 3, 4\}\}.
                    \end{flalign*}
                \end{spreadlines}

                Since $T(n) = \abs{L_n}$, the sequence of $T(n)$ for $1 \le n \le 10$ is
                \[ \left(T(n)\right)_{n \le 10} = \{0, 0, 1, 1, 1, 1, 2, 1, 3, 2\}. \]

                \bigskip

            \item
                \begin{proposition} \label{prop:2-b}
                    The sequence $(a_n)_{n \in \mathbb{N}}$ given by $a_n = T(n)$ has the generating function
                    \[ G(x) = \frac{x^3}{(1 - x^2) (1 - x^3) (1 - x^4)}. \]
                \end{proposition}

                \medskip

                \begin{proof}
                    We can intepret $T(n)$ as the number of partitions of $n$ with exactly three parts $a \le b \le c$ such that $a + b > c$.

                    Let $\{a, b, c\}$ be a partition of $n$ where
                    \[ b = a + i, \qquad c = b + j \]
                    for non-negative $i$ and $j$.

                    Note that $c = a + i + j$.
                    Expanding the previous inequality gives
                    \begin{spreadlines}{\baselineskip}
                        \begin{align}
                            && c &< a + b, \nonumber && \\
                            \Longleftrightarrow && (a + i + j) &< a + (a + i), \nonumber \\
                            \Longleftrightarrow && j &< a. \label{eq:j-lt-a}
                        \end{align}
                    \end{spreadlines}
                    This means that $a + b > c$ if and only if $j < a$.

                    To obtain the generating function for $(a_n)_{n \in \mathbb{N}}$ where $a_n = T(n)$, let $m = a + b + c$ and take the sum of $x^m$ for all possible values of $a, b, c$ such that $a + b > c$.
                    For any $n$, the value of $m$ equals $n$ exactly $T(n)$ times because $a, b, c$ sweep over all possible partitions of $m$, for any $m \ge 0$, that describe the lengths of a triangle.

                    To compute this, we can sum over the possible values of $a$, $i = b - a$, and $j = c - b$.
                    Since the parts of $m$ cannot be greater than $m$ and the side lengths of a triangle cannot be zero, the bounds of $a$ and $i$ are $1 \le a \le m$ and $0 \le i \le m - a$.
                    From \cref{eq:j-lt-a}, the bounds of $c$ that satisfy the condition $a + b > c$ are $0 \le j < a$.

                    Note that $m = a + (a + i) + (a + i + j) = 3a + 2i + j$.
                    The generating function of $(a_n)_{n \in \mathbb{N}}$ is thus given by the summation of as $m$ goes to infinity:
                    \begin{spreadlines}{1.5\baselineskip}
                        \begin{align*}
                            G(x) &= \sum_{a = 1}^\infty\sum_{i = 0}^\infty\sum_{j = 0}^{a - 1} x^{3a + 2i + j}, \\
                            &= \sum_{a = 1}^\infty \sum_{i = 0}^\infty \left( x^{3a + 2i} \sum_{j = 0}^{a - 1} x^j \right), \\
                            &= \sum_{a = 1}^\infty \sum_{i = 0}^\infty \left( x^{3a + 2i} \cdot \frac{1 - x^a}{1 - x} \right), \\
                            &= \frac{1}{1 - x} \sum_{a = 1}^\infty \left( x^{3a} (1 - x^a) \sum_{i = 0}^\infty x^{2i} \right), \\
                            &= \frac{1}{1 - x} \sum_{a = 1}^\infty \left( (x^{3a} - x^{4a}) \cdot \frac{1}{1 - x^2} \right), \\
                            &= \frac{1}{(1 - x) (1 - x^2)} \left( \frac{x^3}{1 - x^3} - \frac{x^4}{1 - x^4} \right), \\
                            &= \frac{x^3}{(1 - x) (1 - x^2)} \left( \frac{(1 - x^4) - (x - x^4)}{(1 - x^3) (1 - x^4)} \right), \\
                            &= \frac{x^3}{(1 - x) (1 - x^2) (1 - x^3) (1 - x^4)} (1 - x), \\
                            &= \frac{x^3}{(1 - x^2) (1 - x^3) (1 - x^4)}. && \qedhere
                        \end{align*}
                    \end{spreadlines}
                \end{proof}

                \bigskip

            \item
                \begin{proposition} \label{prop:2-c}
                    The value of $T(n + 9)$ for $n \in \mathbb{N}$ is equal to
                    \[ T(n) - T(n + 2) - T(n + 3) - T(n + 4) + T(n + 5) + T(n + 6) + T(n + 7). \]
                \end{proposition}

                \medskip

                \begin{proof}
                    From \cref{prop:2-b}, the generating function of the sequence $(a_n)_{n \in \mathbb{N}}$ with $a_n = T(n)$ is given by
                    \[ G(x) = \frac{x^3}{(1 - x^2) (1 - x^3) (1 - x^4)}. \]
                    Rearranging, we have
                    \begin{spreadlines}{0.5\baselineskip}
                        \begin{align}
                            x^3 &= (1 - x^2) (1 - x^3) (1 - x^4) \, G(x), \nonumber \\
                            &= (1 - x^2 - x^3 - x^4 + x^5 + x^6 + x^7 - x^9) \, G(x). \label{eq:generating-function}
                        \end{align}
                    \end{spreadlines}

                    Since $G(x) = \sum_{n = 0}^\infty a_n x^n$, the value of $x^a \, G(x)$ for any $a \ge 0$ is
                    \[ x^a \, G(x) = \sum_{n = 0}^\infty T(n) x^{n + a}. \]
                    The coefficient of $x^n$ in this sum is thus $T(n - a)$.

                    We can thence take the coefficient of $x^{n + 9}$ on both sides of \cref{eq:generating-function} to get
                    \begin{spreadlines}{0.5\baselineskip}
                        \begin{align*}
                            && 0 &= T(n + 9) - T(n + 7) - T(n + 6) - T(n + 5) \\
                            && &\quad + T(n + 4) + T(n + 3) + T(n + 2) - T(n), && \\
                            \text{hence} && T(n + 9) &= T(n) - T(n + 2) - T(n + 3) - T(n + 4) \\
                            && &\quad + T(n + 5) + T(n + 6) + T(n + 7). && \qedhere
                        \end{align*}
                    \end{spreadlines}
                \end{proof}

        \end{enumerate}

        \bigskip

    \item On the power set $\mathcal{P}(\mathbb{N})$ of the natural numbers, define the relation $\sim$ as follows:

        \smallskip

        \begin{definition} \label{def:relation}
            For $A, B \in \mathcal{P}(\mathbb{N})$,
            \[ A \sim B \quad \text{if and only if} \quad (A \, \backslash \, B) \cup (B \, \backslash \, A) \; \text{is finite}. \]
        \end{definition}

        \bigskip

        \begin{enumerate}

            \item
                \begin{proposition} \label{prop:3-a}
                    The relation $\sim$ is an equivalence relation.
                \end{proposition}

                \medskip

                \begin{proof}
                    We shall prove that the relation $\sim$ is reflexive, symmetric, and transitive.
                    For the following claims, let $A$, $B$, and $C$ be sets in $\mathcal{P}(\mathbb{N})$.

                    \begin{claim}
                        $A \sim A$ is always true for any set $A$.
                    \end{claim}

                    \medskip

                    \begin{subproof}
                        From \cref{def:relation},
                        \begin{align*}
                            A \sim A &\Longleftrightarrow (A \, \backslash \, A) \cup (A \, \backslash \, A) \; \text{is finite}, \\
                            &\Longleftrightarrow \varnothing \; \text{is finite}.
                        \end{align*}

                        This is always a true statement, hence $A \sim A$ is always true.
                    \end{subproof}

                    \medskip

                    \begin{claim}
                        $A \sim B$ if and only if $B \sim A$.
                    \end{claim}

                    \medskip

                    \begin{subproof}
                        Since the union of sets is symmetric, we have
                        \[ (A \, \backslash \, B) \cup (B \, \backslash \, A) \Longleftrightarrow (B \, \backslash \, A) \cup (A \, \backslash \, B). \]

                        The LHS is equal to $A \sim B$ and the RHS is equal to $B \sim A$, hence the claim.
                    \end{subproof}

                    \medskip

                    \begin{claim}
                        If $A \sim B$ and $B \sim C$, then $A \sim C$.
                    \end{claim}

                    \begin{subproof}
                        We wil first show that if $U, V, W \in \mathcal{P}(\mathbb{N})$, then
                        \begin{equation} \label{eq:union}
                            (U \, \backslash \, V) \cup (V \, \backslash \, W) = U \, \backslash \, W.
                        \end{equation}

                        Consider an element $x \in (U \, \backslash \, W)$.
                        Then, \cref{eq:union} is equivalent to the statement
                        \[ x \in U \; \text{and} \; x \notin V \quad \text{or} \quad x \in V \; \text{and} \; x \notin W. \]

                        But we already have that $x \in U$ and $x \notin W$.
                        For either of the above conditions to be true, either $x \in V$ or $x \notin V$.
                        This is always true, hence \cref{eq:union} is always true.

                        Suppose that $A \sim B$ and $B \sim C$. Then,
                        \begin{flalign*}
                            && & (A \, \backslash \, B) \cup (B \, \backslash \, A) \; \text{is finite} && \\
                            \text{and} && & (B \, \backslash \, C) \cup (C \, \backslash \, B) \; \text{is finite}.
                        \end{flalign*}
                        The union of two sets is finite if and only if both sets are finite, hence
                        \[ \left((A \, \backslash \, B) \cup (B \, \backslash \, A)\right) \cup \left((B \, \backslash \, C) \cup (C \, \backslash \, B)\right) \; \text{is finite}. \]
                        Since the union is associative, we can rearrange to get
                        \[ \left((A \, \backslash \, B) \cup (B \, \backslash \, C)\right) \cup \left((C \, \backslash \, B) \cup (B \, \backslash \, A)\right) \; \text{is finite}, \]
                        and by using \cref{eq:union},
                        \[ (A \, \backslash \, C) \cup (C \, \backslash \, A) \; \text{is finite}. \]

                        This is equal to $A \sim C$, hence
                        \[ A \sim B \; \text{and} \; B \sim C \implies A \sim C. \]
                    \end{subproof}

                    \medskip

                    We have thus proven that $\sim$ is reflexive, symmetric, and transitive, hence $\sim$ is an equivalence relation.
                \end{proof}

                \bigskip

            \item Let $X$ be an arbitrary subset of $\mathbb{N}$.

                The statement $X \sim \varnothing$ simplifies to
                \[ (X \, \backslash \, \varnothing) \cup (\varnothing \, \backslash \, X) \; \text{is finite} \quad \Longleftrightarrow \quad X \; \text{is finite}. \]
                Therefore, the equivalence class of $\varnothing$ is the set of finite subsets of $\mathbb{N}$.

                Now consider the statement $X \sim \mathbb{N}$. This simplifies to
                \[ (X \, \backslash \, \mathbb{N}) \cup (\mathbb{N} \, \backslash \, X) \; \text{is finite} \quad \Longleftrightarrow \quad (\mathbb{N} \, \backslash \, X) \; \text{is finite}. \]

                Since $\mathbb{N}$ is infinite, $(N \, \backslash \, X)$ is finite if and only if $X = (\mathbb{N} \, \backslash \, Y)$ for some finite subset $Y$ of $\mathbb{N}$.
                This means that the equivalence class of $\mathbb{N}$ is
                \[ [\mathbb{N}] = \{(\mathbb{N} \, \backslash \, Y) \; \big| \; Y \subseteq \mathbb{N} \; \text{and} \; Y \; \text{is finite}\}. \]

                \bigskip

            \item
                \begin{proposition} \label{prop:3-c}
                    There are an infinite number of equivalence classes of $\sim$ on $\mathcal{P}(\mathbb{N})$.
                \end{proposition}

                \begin{proof}
                    Define $S_n$ to be the set $\{k n \; \big| \; k \ge 0\}$ for $n \in \mathbb{N}$.
                    Note that $S_n$ is necessarily a subset of $\mathbb{N}$.

                    Consider the statement $S_1 \sim S_2$.
                    We have $S_1 \, \backslash \, S_2 = \{1, 3, 5, \ldots\}$ which is infinite, so $S_1 \sim S_2$ is false.
                    This means that $S_1$ and $S_2$ belong to different equivalence classes.
                    With the same logic, we can determine that $S_1$, $S_2$, and $S_3$ belong to different equivalence classes.
                    Indeed, for any choice of $l, k \in \mathbb{N}$ with $l < k$, an equivalence class contains both $S_m$ and $S_n$ if and only $S_k \, \backslash \, S_l$ is finite.
                    This is never true, since there is always an infinite number of elements in $S_k$ that are not in $S_l$.
                    Hence for each $n \in \mathbb{N}$, the set $S_n$ describes its own equivalence class.

                    In this sense, there is a one-to-one mapping from $\mathbb{N}$ to the equivalence classes of $\sim$ on $\mathcal{P}(\mathbb{N})$. Since $\mathbb{N}$ is an infinite set, it follows that there are an infinite number of equivalence classes.
                \end{proof}

        \end{enumerate}

\end{enumerate}

\pagebreak

\section*{Feedback}

For \cref{prop:1-c} (question 1.c), I found it particularly difficult to contruct a graph with the required radius and diameter that was simple enough to generalise mathematically. Eventually I succeeded, but I wonder if there is an even easier approach to this problem. I would also like feedback on the rigour of my mathematics with regard to the proof with this graph. Did I make any assumptions that require further explanation?

The way that I approached \cref{prop:3-c} (question 3.c) also seemed to assume too much of the reader. Again, I would like feedback on whether there is a more straightforward, or perhaps more elegant way of arguing my point. Supposing that I'm correct, of course.

\pagebreak

\section*{Acknowledgement}

I did not use AI or any external resources to complete this submission.
